% !TeX root = ../main.tex

\section{AQM Problem Set 2 -- 4}

\begin{problem}[Coherent-state path integral]
   Let's recall that a bosonic coherent state is defined as
  \[
    \ket|\alpha> = D(\alpha) \ket|0>, \quad
    D(\alpha) \equiv \upe^{\alpha\hat a^\dagger-\alpha^*\hat a},
  \]
  where $D(\alpha)$ is the unitary displacement operator. The
  set $\{\ket|\alpha>, \alpha \in \mathbb C\}$ forms an overcomplete basis of
  the Hilbert space
  \[
    \int \frac{\d^2\alpha}{\pi} \ketbra|\alpha><\alpha| = \identity,
  \]
  where $\d^2\alpha = \d\Re\alpha \d\Im\alpha
    = \frac{1}{2\iu}\d\bar\alpha \wedge \d\alpha$
  ($\alpha = \Re\alpha + \iu\Im\alpha$, $\bar\alpha \equiv \alpha^*$), and
  \[
    \braket<\alpha|\beta>
  = \upe^{-\frac12(\bar\alpha \alpha+\bar\beta\beta) + \bar\alpha \beta}
  \]
  We now continue to explore the use of coherent states to do path integral.
  \begin{enumext*}[label = (\alph*)]
    \item Consider a Hamiltonian $\hat H(\hat a^\dagger, \hat a; t)$ which is
    normal ordered (namely, no $\hat a^\dagger$ appears to the right of $\hat a$
    in a product), and the propagator
    \[
      \braket*\Big<\psi_f|\hat U(t_f, t_i)|\psi_i>
    = \braket*\Big<\psi_f|\mathcal T \exp\ab\Big[
        -\iu\int_{t_i}^{t_f} \d t \hat H(\hat a^\dagger, \hat a; t)]|\psi_i>
    \]
    from an initial state $\ket|\psi_i>$ to a final state $\ket|\psi_f>$, where
    $\mathcal T$ is the time-ordering operator. Derive the coherent-state path
    integral representation of the propagator by inserting resolutions of
    identity in terms of $\{\ket|\alpha>\}$
    (Hint: $\braket<\alpha|\hat H(\hat a^\dagger, \hat a; t)|\beta>
      = H(\bar\alpha, \beta; t)\braket<\alpha|\beta>$ with the normal ordering).
    Optional: compare the result with the canonical phase-space path
    integral representation.
    \item* (optional) Consider, again, the specific case of a driven
    quantum harmonic oscillator with Hamiltonian
    \[
      \hat H = \frac{\hat p^2}{2m} + \frac12 m\omega^2 \hat x^2
             + \mathcal E(t) \hat x.
    \]
    Compute the transition amplitude (as a function of time) from the
    unperturbed ground state $\ket|0>$ to an unperturbed excited state
    $\ket|n>$ by using the coherent-state path integral representation of the
    propagator derived above. By treating the driving
    term $\mathcal E(t) \hat x$ as a source term and the path integral as a
    generating functional $Z|\mathcal E(t)|$, compute the two-point correlation
    function $\braket<x(t_1)x(t_2)>$ and the four-point correlation function
    $\braket<x(t_1)x(t_2)x(t_3)x(t_4)>$.
  \end{enumext*}
\end{problem}
\begin{solution}

\end{solution}

\begin{problem}[Spin path integral]
  Let's now move on to do path integral with a single spin-$1/2$.
  \begin{enumext}
    \item Define spin coherent state as
    \[
      \ket|\alpha> = D(\alpha) \ket|\uparrow>, \quad
      D(\alpha) \equiv \upe^{\alpha\hat J_--\alpha^*\hat J_+},
    \]
    where $\hat \sigma_z\ket|\uparrow> = \ket|\uparrow>$
    (similarly, $\hat \sigma_z\ket|\downarrow> - -\ket|\downarrow>$)
    and $\hat J_\pm = (\hat \sigma_x \pm \iu\hat \sigma_y)/2$ with
    $\hat\sigma_{x,y,z}$ the Pauli operators. Show that
    \[
      \ket|\alpha> = \cos|\alpha| \ket|\uparrow>
    + \sin|\alpha| \upe^{\iu\varphi_\alpha} \ket|\downarrow>\quad
      (\alpha \equiv |\alpha| \upe^{\iu\varphi_\alpha}).
    \]
    From now on, we will relabel $\alpha$ by the unit normal vector $\bm n$ on
    the Bloch sphere such that the spin coherent state becomes
    \begin{multline*}
      \ket|\bm n = (\sin\theta\cos\varphi, \sin\theta\sin\varphi, \cos\theta)>
    \equiv \ket*|\alpha = \frac\theta2\upe^{\iu\varphi}>\\
    = \cos\frac\theta2\ket|\uparrow>
    + \sin\frac\theta2\upe^{\iu\varphi} \ket|\downarrow>, \quad
      (\theta \in [0, \pi],\ \varphi \in [0, 2\pi]).
    \end{multline*}
    Show that
    \[
      \braket<\bm n|\bm B \cdot \hat{\bm \sigma}|\bm n> = \bm B \cdot \bm n,
    \]
    where $\hat{\bm \sigma} = (\hat \sigma_x, \hat \sigma_y, \hat \sigma_z)$
    and $\bm B$ is an arbitrary real vector.
    \item Show that
    \[
      \int \frac{\d^2\bm n}{2\pi} \ketbra|n><n| = \identity,
    \]
    where $\d^2\bm n = \sin\theta\d\theta\d\varphi$ is the area element on the
    unit sphere. Further show that
    \[
      \braket<\bm n|\bm n'> = \ab(\frac{1 + \bm n \cdot \bm n'}{2})^{\frac12}
      \upe^{\iu\Omega(\hat z, \bm n, \bm n')/2},
    \]
    where $\Omega(\hat z, \bm n, \bm n')$ is the solid angle subtended
    by $(\hat z, \bm n, \bm n')$ such that
    $\tan(\Omega(\hat z, \bm n, \bm n')/2) = \hat z \cdot (\bm n \times \bm n')/
      (1 + \hat z \cdot \bm n + \bm n \cdot \bm n' + \bm n' \cdot \hat z)$.
    \item Consider Hamiltonian $\hat H = \bm B \cdot \hat{\bm \sigma}$, derive
    the following path integral representation of the partition function
    $Z = \Tr\upe^{-\beta\hat H}$ by using the spin coherent states
    \[
      Z = \exp\ab\{-\frac\iu2 \Omega[\bm n(\tau)]
        - \int_0^\beta \d\tau \bm B \cdot \bm n(\tau)\},
    \]
    where $\Omega[\bm n(\tau)] = \int_0^\beta \d \tau[1 - \cos\theta(\tau)]
    \dot\varphi(\tau)$ is the solid angle subtended by the closed
    path $\bm n(\tau)$ with $\bm n(0) = \bm n(\beta)$. Optional: From this
    result, find the classical equation of motion for $\bm n$ of real time $t$.
  \end{enumext}
\end{problem}
\begin{solution}
  
\end{solution}

\begin{problem}[Discrete translational symmetry]
  Consider the following Hamiltonian for a 1D system
  \[
    \hat H = \frac{\hbar^2}{2m}\hat k^2 + V\cos(2\pi\hat x/\lambda),
  \]
  where $\lambda$ is a positive real number.
  \begin{enumext*}[label = (\alph*)]
    \item Compute the commutator $[\hat k, \hat H]$. Is the eigenvalue
    of $\hat k$ a good quantum number in this problem?
    \item Define $\hat T_\lambda = \upe^{-\iu\lambda\hat k}$, show that
    $\hat T_\lambda\ket|x>$ is an eigenstate of $\hat x$. What is the physical
    meaning of $\hat T_\lambda$ as a transformation of quantum states? Compute
    the comtator $[\hat T_\lambda, \hat H]$. Is the eigenvalue
    of $\hat T_\lambda$ a good quantum number in this problem?
    \item Find all engenstates of $\hat T_\lambda$ with eigenvalue $1$ and $-1$,
    respectively. Write down the projected Hamiltonian in these two
    fixed $\hat T_\lambda$-eigenvalue sub-Hilbert-spaces, respectively.
    \item* (optional) We will call the $\hat T_\lambda$-eigenvalue $1$
    (resp., $-1$) subspace the $k = 0$ (resp., $k = \pi$) subspace. To the
    second order in $V$ as perturbation, compute the approximate energy and the
    corresponding wave-function of the lowest-energy eigenstate of the
    Hamiltonian in the $k = 0$ space, and compute the first excitation gap
    (i.e. the energy difference between the lowest and the second-lowest energy
    eigenstates of the Hamiltonian) in the $k = \pi$ space.
  \end{enumext*}
\end{problem}
\begin{solution}
  
\end{solution}

\begin{problem}[Representation of $D_3$]
  Let's extend the example used in our classes by including two
  flavors ($s = \pm$) on each vertices of the equilateral triangle, such that
  the Hilbert space becomes $V \equiv \operatorname{span}
    (\{\ket|i = 9,1,2; s = \pm>\})$. The action of the two generators of the
  dihedral group $D_3 = \braket<R,\rho|R^3 = \rho^2 = (\rho R)^2 = e>$ on this
  Hilbert space is given by
  \[
    \hat R\ket|i, s> = \ket|i + 1, s>, \quad
    \hat \rho \ket|i, s> = \ket|-i, -s>.
  \]
  Here, $\hat R$ and $\hat \rho$ are unitary operators, and integers such as
  $i - 1$ and $-i$ are understood to be modulo $3$.
  \begin{enumext}
    \item Check that as operators, $\hat R^3 = \hat\rho^2
    = (\hat \rho\hat R)^2 = \identity$, such that they indeed generate a
    representation of $D_3$.
    \item Derive the most general form of the Hamiltonian $\hat H$ on $V$ which
    is invariant under the action of $D_3$, i.e., $[\hat H, \hat R]
    = [\hat H, \rho] = 0$. How many independent (real) parameters does it have?
    \item Consider the following Hamiltonian
    \[
      \hat H = t(\hat R_+ \oplus \hat R_-^2)
             + t^*(\hat R_+^2 \oplus \hat R_-)
             + s(\hat \rho_0 + \hat \rho_1 + \hat \rho_2),
    \]
    where $\hat R_{s=\pm}$ is the restriction of $\hat R$ on the
    subspace $V_s \equiv \operatorname{span}(\{\ket|i = 0,1,2;s>\})$
    (such that $\hat R = \hat R_+ \oplus \hat R_-$);
    $\hat \rho_{i=0,1,2} \equiv \hat\rho \hat R^i$ is the reflection with
    respect to the $i$-th vertex; $t$ is a complex parameter and $s$ is a real
    parameter. Check that this Hamiltonian is invariant under the action
    of $D_3$. Further check that the three parts of the Hamiltonian: $\hat H_R
    \equiv t(\hat R_+ \oplus \hat R_-^2)$, $\hat H_R^\dagger
    = t^*(\hat R_+^2 \oplus \hat R_-)$
    and $\hat H_\rho \equiv s(\hat \rho_0 + \hat \rho_1 + \hat \rho_2)$ all
    commute with each other: $[\hat H_R, \hat H_R^\dagger] = 0$,
    $[\hat H_R, \hat H_\rho] = 0$, and $[\hat H_R^\dagger, \hat H_\rho] = 0$.
    Use this preparation to solve the eigensystem of $\hat H$ exactly. Analyze
    the result based on the irreps of $D_3$.
  \end{enumext}
\end{problem}
\begin{solution}
  
\end{solution}

\begin{problem}[Angular momentum]
  Let $\bm J = (J_x, J_y, J_z)$ be an angular momentum operator, where $J_x$,
  $J_y$, $J_z$, and hence $J_\pm \equiv J_x \pm \iu J_y$, satisfy the usual
  angular-momentum commutation relations.
  \begin{enumext*}[label = (\alph*)]
    \item Show that $\bm J^2 = J_z^2 + J_+J_- - \hbar J_z$.
    \item Starting from $\bm J^2\ket|j, m> = \hbar^2 j(j + 1) \ket|j, m>$
    and $J_z\ket|j, m> = \hbar m \ket|j, m>$, and making use of the above
    equation, derive the coefficient $c_m$ defined
    in $J_-\ket|j, m> = \hbar c_m\ket|j, m - 1>$.
    \item Based on the above result, write down the matrix representation of
    $J_x$ and $J_y$ in the $\{\ket|j, m>\}$ basis for the spin-$3/2$
    (namely, $j = 3/2$) case.
    \item Consider the combination of two degrees of freedom: An $l = 1$ orbital
    (i.e., a $p$-orbital) and an $s = 1/2$ spin, write down the new bases
    (in terms of the old ones) for the combined system decomposed according to
    the total angular momentum $j$.
    \item* (optional) The time-reversal operator in the $\{\ket|j, m>\}$ basis
    is given by $\mathcal T = \upe^{\iu\pi J_y/\hbar} K$, where $K$ is the
    complex conjugation operator. Show that $\mathcal T^2 = (-1)^{2j}$.
  \end{enumext*}
\end{problem}
\begin{solution}
  
\end{solution}

\begin{problem}[Two-site Hubbard model]
  Consider the Fock space for two sites and two spin-$1/2$ states on each site,
  generated by $\{c_{1\uparrow}^\dagger, c_{1\downarrow}^\dagger,
  c_{2\uparrow}^\dagger, c_{2\downarrow}^\dagger\}$.
  \begin{enumext*}[label = (\alph*)]
    \item What is the dimension of this Fock space? Write down the basis for
    this Fock space grouped by the total particle number.
    \item Compute the full spectrum of the following Hamiltonian (assuming all
    parameters to be real)
    \[
      H = \sum_{\substack{i=1,2\\s=\uparrow,\downarrow}}
          \epsilon c_{is}^\dagger c_{is}
        - \sum_{s=\uparrow,\downarrow}(tc_{1s}^\dagger c_{2s} + \text{H.c.})
        + \sum_{i=1,2} Uc_{i\uparrow}^\dagger c_{i\uparrow}
                        c_{i\downarrow}^\dagger c_{i\downarrow}.
    \]
    (Hint: Find conserved particle numbers and compute in each
    fixed-particle-number sector separately.)
    \item Further assuming $\epsilon < 0$ and $U \gg (-\epsilon) > t > 0$, find
    the ground state and the first excited state(s) of the Hamiltonian, and
    compute the entangelement entropy between site 1 and site 2 for each of
    these states.
    \item* (optional) In the subspace spanned by
    $(\hat c_{2\uparrow}^\dagger   \hat c_{1\uparrow}^\dagger,
      \hat c_{2\downarrow}^\dagger \hat c_{1\downarrow}^\dagger,
      \hat c_{2\downarrow}^\dagger \hat d_{1\uparrow}^\dagger,
      \hat c_{2\uparrow}^\dagger   \hat c_{1\downarrow}^\dagger) \ket|0>$,
    compute the spectrum and find the ground state (for both $J > 0$ and $J < 0$
    cases) of the following Hamiltonian
    \[
      \hat H_\text{ex} = \sum_{\alpha=x,y,z} J
        (\bm c_1^\dagger \sigma_\alpha \bm c_1)
        (\bm c_2^\dagger \sigma_\alpha \bm c_2),
    \]
    where $\bm c_{i=1,2}^\dagger \equiv (c_{i\uparrow}^\dagger,
    c_{i,\downarrow}^\dagger)$, and $\sigma_{\alpha=x,y,z}$ are Pauli matrices.
    Compare these results with those obtained in part (c).
  \end{enumext*}
\end{problem}
\begin{solution}
  
\end{solution}